%% 03_background.tex
%% Section 3: Institutional Background

\section{Institutional Background}
\label{sec:background}

This section describes the three areas of the coal-banking-sovereign nexus in Indonesia. We first outline the structure and significance of Indonesia's coal sector, then characterise the banking sector's exposure to coal, and finally document the sovereign fiscal dependence on coal revenues. The section concludes by reviewing the evolving regulatory landscape that will shape the pace and cost of the transition.

\subsection{Indonesia's Coal Sector}
\label{subsec:coal_sector}

Indonesia is the world's largest exporter of thermal coal, shipping over 500 million tonnes annually and supplying approximately one-quarter of global seaborne trade \citep{IEA2024}. Total domestic production exceeded 600 million tonnes (Mt) in 2023, ranking Indonesia fourth globally after China, India, and the United States in absolute output, but first among net exporters. Coal contributed approximately 10\% of merchandise export revenue and remains the single largest source of primary energy for domestic power generation, fuelling roughly 60\% of the national electricity mix \citep{IESR2023}.

On coal quality and classification,  Indonesian coal spans a wide range of calorific values, from low-rank sub-bituminous coals with 4,000--5,000 kcal/kg to higher-grade bituminous coals exceeding 6,000 kcal/kg. The majority of production, and of the companies in our sample, falls in the sub-bituminous category, which commands lower prices per tonne but benefits from lower extraction costs due to shallow, surface-mineable deposits. In our dataset, 23 of the 34 companies produce sub-bituminous coal (median calorific value of 4,800 kcal/kg), while the remaining 11 produce bituminous coal (median 6,000 kcal/kg). This distinction matters for stranding analysis because higher-calorie coals command a premium that partially offsets the downward price pressure from the energy transition, while lower-calorie coals face more acute substitution risk from natural gas and renewables.

Regarding the producer, the Indonesian coal sector is characterised by moderate concentration at the top and a long tail of smaller producers. The four largest companies by production volume, PT Adaro Andalan Indonesia (AADI, 65.8 Mt), Bayan Resources (BYAN, 56.9 Mt), Bumi Resources (BUMI, 43.3 Mt), and PT Bukit Asam (PTBA, 43.3 Mt), together account for approximately 44\% of total listed production. PTBA occupies a unique position as the only government-controlled coal miner on the Indonesia Stock Exchange (IDX), serving as both a major revenue source and a policy instrument for the domestic market obligation. Indo Tambangraya Megah (ITMG), a subsidiary of Thailand's Banpu, and Dian Swastatika Sentosa (DSSA) are among the other large-cap producers. At the other end of the spectrum, companies such as Borneo Olah Sarana Sukses (BOSS, 0.3 Mt) and Resources Alam Indonesia (KKGI, 0.5 Mt) operate small concessions with limited remaining mine lives.

 Nearly all Indonesian coal mining is conducted through open-pit (surface) methods, reflecting the shallow geological setting of the major coal basins in South Sumatra, East Kalimantan, and South Kalimantan. The predominance of surface mining implies relatively low extraction costs: the median all-in sustaining cost (AISC) in our sample is approximately \$75 per tonne, though this varies from under \$10 per tonne for ITMA's low-cost operations to over \$169 per tonne for TOBA's higher-cost deposits. Strip ratios (the volume of overburden removed per tonne of coal) are generally low by global standards, though data coverage on this metric is limited in public disclosures. Mine lives are conventionally assumed at 20 years in the absence of detailed reserve data, reflecting the typical duration of Indonesian coal concession contracts (Izin Usaha Pertambangan, or IUP).

Regarding domestic market, a distinctive feature of the Indonesian coal market is the government-mandated Domestic Market Obligation, which requires producers to sell a minimum of 25\% of their planned production to domestic buyers, primarily state electricity utility PT PLN, at a capped price (currently \$70 per tonne for power generation). The DMO compresses margins on domestic sales and effectively creates a cross-subsidy from export earnings to domestic electricity costs. For stranding analysis, the DMO introduces a floor on domestic demand that may slow the pace of coal phase-down but simultaneously reduces per-tonne profitability, making the sector more vulnerable to international price declines \citep{Burke2018, Makhijani2015}.

On export markets and competitive position, Indonesian thermal coal is exported primarily to China, India, Japan, South Korea, and Southeast Asian neighbours. Its competitive advantage rests on geographic proximity to Asian demand centres and low extraction costs, but it faces growing competition from Australian and South African coal at the higher-calorie end and from Russian coal that has been redirected to Asian markets following Western sanctions. Under aggressive decarbonisation scenarios, demand from Japan and South Korea, both of which have announced net-zero targets, is projected to decline sharply by 2035, while Chinese and Indian demand paths remain subject to considerable uncertainty \citep{IEA2023, Mendelevitch2019}.

The recent coal price cycle provides important context for understanding the sector's transition risk exposure. Newcastle benchmark prices collapsed to approximately \$48/tonne in September 2020 as COVID-19 crushed energy demand, then surged to unprecedented levels above \$400/tonne in mid-2022 following Russia's invasion of Ukraine and the ensuing global energy security panic. By 2023--2025, prices normalised to the \$100--140/tonne range as supply adjusted and renewable energy deployment accelerated. This extreme volatility illustrates the commodity price risk to which Indonesian coal producers and their bank lenders are exposed, while the long-run downward trend in prices is broadly consistent with the transition pathways modelled in this paper. Notably, China's unofficial ban on Australian coal imports from late 2020 temporarily redirected demand toward Indonesian suppliers, reinforcing Indonesia's role as the swing exporter in the seaborne thermal coal market.

\subsection{Banking Sector and Coal Exposure}
\label{subsec:banking_exposure}

Indonesia's financial system is bank-dominated. Bank credit finances approximately 34\% of GDP, and listed banks on the IDX collectively hold over \$700 billion in total assets \citep{IMF2024Indonesia}. The sector is highly concentrated: four domestic systemically important banks (D-SIBs), namely Bank Mandiri (BMRI), Bank Rakyat Indonesia (BBRI), Bank Central Asia (BBCA), and Bank Negara Indonesia (BBNI), together account for roughly 55\% of total banking assets and 70\% of total lending. Bank Tabungan Negara (BBTN), while classified as a D-SIB for regulatory purposes, holds negligible coal exposure and is excluded from the coal-banking linkage analysis. Beyond the D-SIBs, Bank CIMB Niaga (BNGA), Bank Permata (BNLI), and Bank Syariah Indonesia (BRIS) are among the banks with material coal-sector lending.

Bank financing to the coal sector flows through three principal channels. First, \textit{project finance} facilities provide long-term funding for mine development, typically secured against mining concessions and future cash flows. Second, \textit{corporate credit facilities}, including revolving working capital lines, trade finance for coal shipments, and term loans for capital expenditure, constitute the bulk of bank-coal exposure in our data. Third, \textit{syndicated loans} involve multiple lenders (often including foreign banks) sharing exposure to large coal projects. Our exposure data, sourced from annual report disclosures, captures the first two channels for IDX-listed banks and attributes syndicated participations where individual bank shares are disclosed.

We attempt to measure the magnitude of exposure. From the exposure matrix constructed in Section~\ref{sec:data}, we identify a bipartite network connecting 13 IDX-listed banks to 26 coal companies, with total identified exposure of approximately \$6.1 billion. Of this, \$3.3 billion is directly attributed to specific bank-company pairs, while \$2.8 billion is classified as unattributed (held by foreign banks, non-bank financial institutions, or syndicated facilities where individual participations are not disclosed). Among IDX-listed banks, Bank Mandiri (BMRI) holds the largest attributed coal exposure at approximately \$1.2 billion, followed by Bank Negara Indonesia (BBNI) at \$975 million and Bank Central Asia (BBCA) at \$441 million.

 The coal exposure is concentrated in a small number of large bank-firm links. The single largest exposure, BMRI's lending to Dian Swastatika Sentosa (DSSA) at approximately \$586 million, represents a material concentration risk by any standard. The five largest bank-coal exposures account for over 50\% of total attributed exposure. This concentration implies that the default or distress of even a small number of coal companies could generate outsized losses for specific banks, a feature we exploit in the stress test of Section~\ref{sec:model2}.

In the financial sector, Indonesian banks are heavily regulated by the Financial Services Authority (Otoritas Jasa Keuangan, OJK), which imposes a minimum capital adequacy ratio (CAR) of 8\%, aligned with Basel III standards. D-SIBs are subject to additional capital buffers, bringing the effective minimum to approximately 10.5\%. As of the latest reporting period, all 38 banks in our sample meet the minimum CAR, with D-SIBs reporting CARs ranging from 19\% to 29\%. However, these ratios are calculated on current risk weights that do not incorporate forward-looking climate transition risk, a gap that motivates our stress test.

Bank lending to the mining sector has expanded rapidly in recent years. According to OJK statistics, mining sector credit outstanding reached Rp~342 trillion (approximately \$21.6 billion) by September 2024, representing 4.5\% of total banking credit and growing at 26.7\% year-on-year. This growth is notable given the sector's historically elevated credit risk: mining sector non-performing loans peaked at 7.7\% in May 2017 during the 2015--2016 commodity price downturn, when coal prices fell below \$50/tonne. The subsequent recovery in coal prices from 2017 onward has improved asset quality and encouraged renewed bank appetite for mining sector lending, but this procyclical pattern raises the concern that banks may be extending credit precisely when transition risk is accumulating.


\subsection{Fiscal Dependence on Coal Revenue}
\label{subsec:fiscal_dependence}

The Indonesian central government derives revenue from the coal sector through three principal channels: production royalties, corporate income tax, and state-owned enterprise dividends. Together, these contributed approximately \$8.5 billion, or 4.5\% of total government revenue, in the 2022 fiscal year, though this share has fluctuated with international coal prices.


 Coal royalties are classified as non-tax state revenue (Penerimaan Negara Bukan Pajak, PNBP) and are levied at rates that depend on coal type and mining stage: 3--5\% for sub-bituminous coal and 5--7\% for bituminous coal. In our dataset, the median royalty rate is 5\%, with higher-calorie producers such as BYAN and ITMG paying 7\%. Royalty receipts are directly linked to coal production volumes and prices, making them among the most commodity-sensitive components of the national budget.

 Coal mining companies are subject to the standard Indonesian corporate income tax rate of 22\%. Given the sector's historical profitability (median EBITDA margins of 25--40\% during the 2021--2023 commodity supercycle), corporate tax contributions from coal have been substantial. However, tax receipts are highly procyclical: a 25\% decline in coal prices (consistent with our 2\textdegree C scenario by 2030) would mechanically reduce taxable income, with the effect amplified by operating leverage given the significant fixed-cost component of mining operations.

The key to the government revenue is dividends from State-owned coal enterprise. PT Bukit Asam (PTBA), as the only publicly listed state-owned coal company, pays dividends to the government through the Ministry of State-Owned Enterprises. With a market capitalisation of approximately \$1.9 billion and a historical dividend yield of roughly 7\%, PTBA contributes an estimated \$130 million annually in dividend income to the state. While this is modest relative to total revenues, it represents a direct and visible fiscal stake in coal sector profitability that complicates political decisions regarding coal phase-down.

 Beyond the central government, coal-producing provinces and regencies in Kalimantan and South Sumatra receive revenue-sharing transfers (Dana Bagi Hasil) from coal royalties, which constitute a major source of sub-national government funding. While this paper focuses on national-level fiscal impacts, the regional dimension amplifies the political economy constraints on the coal transition and underscores the importance of just-transition mechanisms \citep{Heffron2018, JakobSteckel2016}.

\subsection{Regulatory Landscape and Climate Finance Policy}
\label{subsec:regulatory_landscape}

The institutional architecture governing Indonesia's climate finance response has evolved rapidly in recent years, though a significant gap remains between policy ambition and supervisory practice. Understanding this landscape is essential for contextualising the stress test results and identifying where regulatory action can most effectively reduce the financial risks we document.

\textit{OJK's Sustainable Finance Roadmap.} The Financial Services Authority (OJK) launched Phase I of its Sustainable Finance Roadmap in 2015 and extended it through Phase II (2021--2025), which mandates disclosure of environmental, social, and governance (ESG) risk in annual reports, requires banks with assets above Rp~1 trillion to apply the POJK 51/2017 sustainable finance principles, and establishes green bond and sukuk certification standards. Phase II specifically calls for climate scenario analysis and the incorporation of environmental risk into bank credit risk frameworks. While these requirements represent a significant step, implementation has been uneven: disclosure quality varies markedly across banks, and climate scenario analysis remains aspirational rather than prescribed in standardised form. Our stress test provides the kind of quantitative, scenario-based benchmark that OJK's Phase II framework anticipates but has not yet operationalised.

\textit{Indonesia's NDC and coal phase-down commitments.} Indonesia's updated Nationally Determined Contribution (NDC), submitted in 2022, commits to a 31.89\% unconditional reduction in greenhouse gas emissions below the BAU trajectory by 2030, rising to 43.2\% with international support. The NDC includes an explicit coal phase-down target: no new coal power plant construction after 2023 (with exceptions for those already under construction), and an ambition to retire or convert 2.5~GW of coal capacity by 2025 and 33~GW by 2035. Meeting these targets would directly affect the revenue outlook for coal mining companies and their bank lenders, as domestic demand for coal, currently supported by the DMO regime, would decline under an accelerated retirement schedule.

\textit{JETP implementation.} The Just Energy Transition Partnership concept was first announced at COP26 in Glasgow (November 2021), with Indonesia's partnership formally launched at COP27 in Sharm el-Sheikh (November 2022). It represents the most ambitious climate finance commitment Indonesia has undertaken. The \$20 billion financing package, comprising \$10 billion from the International Partners Group (G7 plus Norway) and a matched private finance mobilisation target, is channelled through an Energy Transition Mechanism (ETM) designed to accelerate early coal plant retirement and scale up renewable energy capacity. The ETM pipeline, managed jointly by the Asian Development Bank and the Indonesian government, has identified specific coal plants for early retirement and blended finance facilities. However, as of early 2026, implementation has been slower than anticipated: the financing arrangements remain under negotiation, just-transition provisions for affected workers and communities are not fully specified, and the role of Indonesian banks in refinancing retired coal assets has yet to be formalised. Our analysis of banking sector vulnerability directly informs the design of these JETP financing instruments: banks with elevated coal exposure cannot be expected to simultaneously absorb coal credit losses and scale up renewable finance without targeted regulatory support.

\textit{Bank Indonesia and the green taxonomy.} Bank Indonesia (BI) has taken complementary steps at the monetary policy level. The Bank Indonesia Sustainable Finance Taxonomy (Green Taxonomy), published in 2022 and aligned with ASEAN Taxonomy for Sustainable Finance, classifies economic activities by their environmental impact and provides a framework for green bond eligibility and preferential refinancing treatment. BI's 2023 circular (SEBI No. 14/2023) encourages banks to disclose their exposure to taxonomy-classified activities, creating an indirect incentive to reduce brown lending. BI has also piloted a green repo facility, offering preferential rates on repurchase agreements collateralised by green bonds, though the scale remains limited relative to the banking system's total assets.

The overall picture is one of a regulatory architecture that has the right ambition but lacks the empirical underpinning to calibrate its specific instruments. Capital add-ons for coal-exposed lending require quantitative estimates of the probability of default and loss given default under climate scenarios — precisely the inputs our Model 2 provides. Concentration limits on coal sector lending require a clear picture of which banks hold material exposures and how these exposures interact through the network — precisely the mapping our bipartite network supplies. The gap between policy ambition and supervisory practice is thus not merely a matter of regulatory will; it reflects a genuine information deficit that this paper helps to close.



